\chapter{Simetria}\label{ch:g4:symmetry}

Borboletas, rostos, flocos de neve: a natureza adora figuras cujas
duas \omterm{def:g3:division:half}{metades} coincidem quando dobradas. Este capítulo explora a
simetria com dobras e papel quadriculado --- as construções com o
esquadro esperam o \cref{ch:g6:symmetry}.

\section{Dobrar e sobrepor}

\begin{definition}[Eixo de simetria]\label{def:g4:symmetry:def}
Uma reta é um \emph{eixo de simetria}\index{eixo de simetria} de uma
figura quando, dobrando a figura ao longo dessa reta, as duas \omterm{def:g3:division:half}{metades}
coincidem \emph{exatamente}. Uma figura pode não ter eixo nenhum, ter
um eixo, ou ter vários.
\end{definition}

\begin{omfigure}
\begin{tikzpicture}[scale=0.6]
  % forma simétrica: triângulo isósceles
  \draw[thick, fill=omDef!12] (0,0) -- (2.4,0) -- (1.2,2.6) -- cycle;
  \draw[omThm, dashed, thick] (1.2,-0.5) -- (1.2,3.0);
  \node[below, font=\small] at (1.2,-0.7) {um eixo};
  % retângulo: dois eixos
  \begin{scope}[xshift=5.4cm]
  \draw[thick, fill=omDef!12] (0,0.3) rectangle (3.2,2.3);
  \draw[omThm, dashed, thick] (1.6,-0.4) -- (1.6,3.0);
  \draw[omThm, dashed, thick] (-0.5,1.3) -- (3.7,1.3);
  \node[below, font=\small] at (1.6,-0.7) {dois eixos};
  \end{scope}
  % escaleno: nenhum
  \begin{scope}[xshift=11.2cm]
  \draw[thick, fill=omDef!12] (0,0) -- (3.0,0.4) -- (0.8,2.2) -- cycle;
  \node[below, font=\small] at (1.4,-0.7) {nenhum eixo};
  \end{scope}
\end{tikzpicture}

{\small Dobre pela linha tracejada: as \omterm{def:g3:division:half}{metades} coincidem. No último
triângulo, nenhuma dobra funciona --- ele não tem \omterm{def:g4:symmetry:def}{eixo de simetria}.}
\end{omfigure}

\begin{example}[Testar com papel vegetal]\label{ex:g4:symmetry:trace}
Decalque a figura, vire o papel vegetal ao longo do eixo candidato e
recoloque-o: se o decalque cobrir a figura exatamente, o eixo é
verdadeiro. Os olhos sozinhos se enganam com frequência --- a
diagonal de um retângulo \emph{parece} um eixo, mas o teste da dobra
diz que não (experimente!).
\end{example}

\section{Completar figuras na malha}

\begin{method}[Simétrico de uma figura em relação a uma linha da malha]\label{met:g4:symmetry:grid}
O eixo é uma linha vertical (ou horizontal) da malha.
\begin{enumerate}
  \item Tome cada \omterm{def:g4:geometry:polygon}{vértice} da figura, um de cada vez;
  \item conte a distância dele até o eixo, em quadradinhos;
  \item ponha o \omterm{def:g4:geometry:polygon}{vértice} imagem à \emph{mesma} distância do
        \emph{outro} lado, na mesma linha (ou coluna);
  \item ligue os \omterm{def:g4:geometry:polygon}{vértices} imagem na mesma ordem e confira: a imagem é
        a gêmea espelhada da figura.
\end{enumerate}
\end{method}

\begin{omfigure}
\begin{tikzpicture}[scale=0.55]
  \draw[gray!40, thin] (0,0) grid (12,5);
  \draw[omThm, very thick] (6,-0.4) -- (6,5.4);
  \draw[thick, omDef, fill=omDef!15]
    (1,1) -- (5,1) -- (5,2) -- (3,2) -- (3,4) -- (1,4) -- cycle;
  \draw[thick, omMeth, fill=omMeth!15]
    (11,1) -- (7,1) -- (7,2) -- (9,2) -- (9,4) -- (11,4) -- cycle;
  \draw[dashed] (5,2) -- (7,2);
  \draw[dashed] (3,4) -- (9,4);
\end{tikzpicture}

{\small Completando em relação ao eixo vermelho: cada \omterm{def:g4:geometry:polygon}{vértice} salta
para a mesma distância do outro lado. \omterm{def:g4:geometry:polygon}{Vértices} que tocam o eixo
ficariam parados.}
\end{omfigure}

\begin{example}[Mesma forma, mesmo tamanho, invertida]\label{ex:g4:symmetry:properties}
A imagem tem exatamente os mesmos comprimentos e os mesmos ângulos
que o original: a simetria copia a figura --- mas a vira do avesso,
como uma mão esquerda e uma mão direita. Se a letra original é F, a
imagem espelhada é um F ao contrário, e não outro F.
\end{example}

\section{Eixos das formas usuais}

\begin{example}[Contar eixos]\label{ex:g4:symmetry:shapes}
Dobrando:
\begin{itemize}
  \item um quadrado: $4$ eixos (dois pelos meios dos lados opostos,
        dois pelas diagonais);
  \item um retângulo: $2$ eixos (só pelos meios dos lados opostos!);
  \item um triângulo isósceles: $1$; um triângulo equilátero: $3$;
  \item um círculo: toda reta que passa pelo centro --- mais eixos do
        que se consegue contar.
\end{itemize}
\end{example}

\begin{example}[Simetria no alfabeto]\label{ex:g4:symmetry:letters}
A, M, T, U, V têm eixo vertical; B, C, D, E têm eixo horizontal; H,
I, O, X têm os dois. Palavras feitas das letras certas também podem
ser simétricas: OVO tem eixo vertical; OXO funciona nos dois
sentidos.
\end{example}

\section{Exercícios}

\begin{exercise}[$\star$]\label{exo:g4:symmetry:1}
Decalque e dobre: quais destas têm \omterm{def:g4:symmetry:def}{eixo de simetria} --- um quadrado;
um retângulo (não quadrado); um triângulo escaleno (todos os lados
diferentes); um círculo?
\end{exercise}

\begin{exercise}[$\star$]\label{exo:g4:symmetry:2}
Trace os \omterm{def:g4:symmetry:def}{eixos de simetria} de: um quadrado; um retângulo; um
triângulo equilátero. Quantos para cada um?
\end{exercise}

\begin{exercise}[$\star$]\label{exo:g4:symmetry:3}
Use o teste do papel vegetal do \cref{ex:g4:symmetry:trace} para
mostrar que a diagonal de um retângulo $6 \times 4$ \emph{não} é um
\omterm{def:g4:symmetry:def}{eixo de simetria}.
\end{exercise}

\begin{exercise}[$\star$]\label{exo:g4:symmetry:4}
No papel quadriculado, trace um eixo vertical e a letra L (três
quadradinhos de altura, dois de largura) a $2$ quadradinhos do eixo.
Construa a imagem espelhada dela.
\end{exercise}

\begin{exercise}[$\star$]\label{exo:g4:symmetry:5}
No papel quadriculado, trace um eixo horizontal e uma bandeirinha (um
triângulo $1 \times 1$ num mastro de $3$ quadradinhos) acima dele.
Complete a figura por simetria abaixo do eixo.
\end{exercise}

\begin{exercise}[$\star$]\label{exo:g4:symmetry:6}
Uma figura toca o seu \omterm{def:g4:symmetry:def}{eixo de simetria} num ponto $P$. Onde fica a
imagem de $P$? Explique com o desenho da dobra.
\end{exercise}

\begin{exercise}[$\star$]\label{exo:g4:symmetry:7}
Classifique as letras maiúsculas A, B, E, F, H, N, O, T: eixo
vertical? eixo horizontal? os dois? nenhum?
\end{exercise}

\begin{exercise}[$\star$]\label{exo:g4:symmetry:8}
Um triângulo tem lados de $5$ cm, $5$ cm e $3$ cm. Ele tem \omterm{def:g4:symmetry:def}{eixo de
simetria}? Quais lados a dobra troca de lugar?
\end{exercise}

\begin{exercise}[$\star\star$]\label{exo:g4:symmetry:9}
No papel quadriculado, desenhe uma figura sua que tenha
\emph{exatamente dois} \omterm{def:g4:symmetry:def}{eixos de simetria}, e trace os dois eixos.
(Dica: pense num retângulo, ou numa cruz.)
\end{exercise}

\begin{exercise}[$\star\star$]\label{exo:g4:symmetry:10}
É dada a \omterm{def:g3:division:half}{metade} de um desenho simétrico: a \omterm{def:g3:division:half}{metade} esquerda de uma
casa (uma parede quadrada e meio telhado triangular), com um eixo
vertical ao longo da borda direita dela. Descreva ou desenhe a casa
completa. Como os comprimentos da \omterm{def:g3:division:half}{metade} direita completada se
comparam aos da \omterm{def:g3:division:half}{metade} esquerda?
\end{exercise}

\begin{exercise}[$\star\star$]\label{exo:g4:symmetry:11}
Aya afirma: ``todo triângulo com dois lados iguais tem \omterm{def:g4:symmetry:def}{eixo de
simetria}, e todo triângulo com \omterm{def:g4:symmetry:def}{eixo de simetria} tem dois lados
iguais.'' Teste a dupla afirmação dela em desenhos: um triângulo
isósceles, um equilátero e um escaleno. A dobra dá razão a ela?
\end{exercise}
